\subsection{Equations with Two Radicals} 
We can also solve equations with two radicals.
\begin{procedure}{Solving an Equation with Two Radicals}
\begin{enumerate*}
\item Isolate and eliminate one radical.
\item Simplify as much as possible. If we've worked correctly, there should now be only one radical in the equation.
\item Isolate and eliminate the remaining radical.
\item Check for extraneous solutions, if necessary.
\end{enumerate*}
\end{procedure}
In theory, we could solve in this way with any number of radicals; but the algebra quickly becomes unweildy.
\begin{Example}
Solve $\sqrt{2x+1}-\sqrt{x+4}=1$.
\begin{lpic}{}{15}
\foreach \y/\lbl in {2/{Eliminate $\sqrt{2x+1}$},
										4/{Simplify},
										7/{Eliminate $\sqrt{x+4}$},
										10/{Finish solving},
										14/{Check}}
	\regtextright{5.75}{-\y}{\lbl};
\draw (6,0) -- (6,-15);
\end{lpic}
\begin{Note}Since we took an \makeblank[1in] power to eliminate the roots, we \emph{must} check for extraneous solutions.\end{Note}
\end{Example}